% コマ11: 式走査の網羅 — コマ10 が素通りした枝を、コマ11 の走査が全部たどる
% printf("%s\n", n < 2 ? "then side" : "else side"); の AST
\documentclass[border=8pt]{standalone}
% 講義資料の図(materials/sessions/figures/*.tex)共通プリアンブル
% 各図の .tex から % 講義資料の図(materials/sessions/figures/*.tex)共通プリアンブル
% 各図の .tex から % 講義資料の図(materials/sessions/figures/*.tex)共通プリアンブル
% 各図の .tex から \input{figure-preamble} で読み込む。
% 配色・フォントは latex/session-template.tex と揃える。

% 日本語(LuaLaTeX + luatexja)。等幅セル内の日本語も和文フォントで出す
\usepackage{luatexja}
\usepackage{luatexja-fontspec}
\setmainfont{Noto Sans CJK JP}
\setmainjfont{Noto Sans CJK JP}
\setmonofont{Inconsolata}[Scale=MatchLowercase]
\setmonojfont{Noto Sans CJK JP}

\usepackage{xcolor}
\definecolor{TcAccentDark}{HTML}{583B66}
\definecolor{TcAccentLight}{HTML}{EEE6F2}
\definecolor{TcAccentLighter}{HTML}{F7F3F9}
\definecolor{TcEmph}{HTML}{E64B17}
\definecolor{TcText}{HTML}{262626}
\definecolor{TcGrayText}{HTML}{737373}
\definecolor{TcGrayBg}{HTML}{F2F2F2}

\usepackage{tikz}
\usetikzlibrary{arrows.meta,positioning,calc,decorations.pathreplacing,fit,backgrounds}

\tikzset{
  % テキスト
  every node/.style={text=TcText},
  annot/.style={font=\small, text=TcText, align=left},
  gannot/.style={font=\small, text=TcGrayText, align=center},
  addr/.style={font=\ttfamily\small, text=TcGrayText},
  % メモリセル
  memcell/.style={draw=TcAccentDark, line width=0.6pt, fill=white,
                  minimum height=8.5mm, minimum width=40mm,
                  font=\ttfamily, inner sep=2pt, outer sep=0pt},
  memcell gray/.style={memcell, fill=TcGrayBg},
  memcell hi/.style={memcell, fill=TcAccentLighter},
  memcell dashed/.style={memcell, dashed, fill=none, text=TcGrayText},
  % 領域(セクション図など)
  region/.style={draw=TcAccentDark, line width=0.6pt, fill=white,
                 minimum width=48mm, align=center, inner sep=3pt, outer sep=0pt},
  % 制御フロー図のボックス
  cfgbox/.style={draw=TcAccentDark, line width=0.6pt, fill=white, rounded corners=1.5mm,
                 minimum width=28mm, minimum height=8mm, font=\small, align=center,
                 inner sep=2pt},
  % 矢印
  ptr/.style={-{Stealth[length=2.8mm]}, line width=1.1pt, color=TcEmph},
  flow/.style={-{Stealth[length=2.4mm]}, line width=0.8pt, color=TcAccentDark},
  flow back/.style={flow, dashed},
  regptr/.style={-{Stealth[length=2.8mm]}, line width=1.1pt, color=TcAccentDark},
  % 補助
  bracestyle/.style={decorate, decoration={brace, amplitude=4pt}, line width=0.8pt,
                     color=TcAccentDark},
}
 で読み込む。
% 配色・フォントは latex/session-template.tex と揃える。

% 日本語(LuaLaTeX + luatexja)。等幅セル内の日本語も和文フォントで出す
\usepackage{luatexja}
\usepackage{luatexja-fontspec}
\setmainfont{Noto Sans CJK JP}
\setmainjfont{Noto Sans CJK JP}
\setmonofont{Inconsolata}[Scale=MatchLowercase]
\setmonojfont{Noto Sans CJK JP}

\usepackage{xcolor}
\definecolor{TcAccentDark}{HTML}{583B66}
\definecolor{TcAccentLight}{HTML}{EEE6F2}
\definecolor{TcAccentLighter}{HTML}{F7F3F9}
\definecolor{TcEmph}{HTML}{E64B17}
\definecolor{TcText}{HTML}{262626}
\definecolor{TcGrayText}{HTML}{737373}
\definecolor{TcGrayBg}{HTML}{F2F2F2}

\usepackage{tikz}
\usetikzlibrary{arrows.meta,positioning,calc,decorations.pathreplacing,fit,backgrounds}

\tikzset{
  % テキスト
  every node/.style={text=TcText},
  annot/.style={font=\small, text=TcText, align=left},
  gannot/.style={font=\small, text=TcGrayText, align=center},
  addr/.style={font=\ttfamily\small, text=TcGrayText},
  % メモリセル
  memcell/.style={draw=TcAccentDark, line width=0.6pt, fill=white,
                  minimum height=8.5mm, minimum width=40mm,
                  font=\ttfamily, inner sep=2pt, outer sep=0pt},
  memcell gray/.style={memcell, fill=TcGrayBg},
  memcell hi/.style={memcell, fill=TcAccentLighter},
  memcell dashed/.style={memcell, dashed, fill=none, text=TcGrayText},
  % 領域(セクション図など)
  region/.style={draw=TcAccentDark, line width=0.6pt, fill=white,
                 minimum width=48mm, align=center, inner sep=3pt, outer sep=0pt},
  % 制御フロー図のボックス
  cfgbox/.style={draw=TcAccentDark, line width=0.6pt, fill=white, rounded corners=1.5mm,
                 minimum width=28mm, minimum height=8mm, font=\small, align=center,
                 inner sep=2pt},
  % 矢印
  ptr/.style={-{Stealth[length=2.8mm]}, line width=1.1pt, color=TcEmph},
  flow/.style={-{Stealth[length=2.4mm]}, line width=0.8pt, color=TcAccentDark},
  flow back/.style={flow, dashed},
  regptr/.style={-{Stealth[length=2.8mm]}, line width=1.1pt, color=TcAccentDark},
  % 補助
  bracestyle/.style={decorate, decoration={brace, amplitude=4pt}, line width=0.8pt,
                     color=TcAccentDark},
}
 で読み込む。
% 配色・フォントは latex/session-template.tex と揃える。

% 日本語(LuaLaTeX + luatexja)。等幅セル内の日本語も和文フォントで出す
\usepackage{luatexja}
\usepackage{luatexja-fontspec}
\setmainfont{Noto Sans CJK JP}
\setmainjfont{Noto Sans CJK JP}
\setmonofont{Inconsolata}[Scale=MatchLowercase]
\setmonojfont{Noto Sans CJK JP}

\usepackage{xcolor}
\definecolor{TcAccentDark}{HTML}{583B66}
\definecolor{TcAccentLight}{HTML}{EEE6F2}
\definecolor{TcAccentLighter}{HTML}{F7F3F9}
\definecolor{TcEmph}{HTML}{E64B17}
\definecolor{TcText}{HTML}{262626}
\definecolor{TcGrayText}{HTML}{737373}
\definecolor{TcGrayBg}{HTML}{F2F2F2}

\usepackage{tikz}
\usetikzlibrary{arrows.meta,positioning,calc,decorations.pathreplacing,fit,backgrounds}

\tikzset{
  % テキスト
  every node/.style={text=TcText},
  annot/.style={font=\small, text=TcText, align=left},
  gannot/.style={font=\small, text=TcGrayText, align=center},
  addr/.style={font=\ttfamily\small, text=TcGrayText},
  % メモリセル
  memcell/.style={draw=TcAccentDark, line width=0.6pt, fill=white,
                  minimum height=8.5mm, minimum width=40mm,
                  font=\ttfamily, inner sep=2pt, outer sep=0pt},
  memcell gray/.style={memcell, fill=TcGrayBg},
  memcell hi/.style={memcell, fill=TcAccentLighter},
  memcell dashed/.style={memcell, dashed, fill=none, text=TcGrayText},
  % 領域(セクション図など)
  region/.style={draw=TcAccentDark, line width=0.6pt, fill=white,
                 minimum width=48mm, align=center, inner sep=3pt, outer sep=0pt},
  % 制御フロー図のボックス
  cfgbox/.style={draw=TcAccentDark, line width=0.6pt, fill=white, rounded corners=1.5mm,
                 minimum width=28mm, minimum height=8mm, font=\small, align=center,
                 inner sep=2pt},
  % 矢印
  ptr/.style={-{Stealth[length=2.8mm]}, line width=1.1pt, color=TcEmph},
  flow/.style={-{Stealth[length=2.4mm]}, line width=0.8pt, color=TcAccentDark},
  flow back/.style={flow, dashed},
  regptr/.style={-{Stealth[length=2.8mm]}, line width=1.1pt, color=TcAccentDark},
  % 補助
  bracestyle/.style={decorate, decoration={brace, amplitude=4pt}, line width=0.8pt,
                     color=TcAccentDark},
}

\begin{document}
\begin{tikzpicture}[node distance=0mm,
    ast/.style={cfgbox, minimum width=30mm, minimum height=7mm, font=\ttfamily\small},
    ast off/.style={ast, dashed, text=TcGrayText, draw=TcGrayText},
    num/.style={annot, font=\small\bfseries, text=TcEmph},
    lab/.style={annot, text=TcEmph, font=\small\ttfamily},
    note/.style={annot, text=TcGrayText, align=left},
    lcline/.style={font=\ttfamily\footnotesize, text=TcText, anchor=west}]

  \node[gannot, anchor=west] at (-58mm, 24mm)
      {{\ttfamily printf("\%s{\textbackslash}n", n < 2 ? "then side" : "else side");} の AST};

  % ---- AST ----
  \node[ast]     (call) at (0mm, 10mm)    {Call printf};
  \node[ast]     (s1)   at (-40mm, -8mm)  {Str "\%s{\textbackslash}n"};
  \node[ast off] (cond) at (24mm, -8mm)   {Cond};
  \node[ast off] (lt)   at (-18mm, -26mm) {Lt};
  \node[ast off] (s2)   at (26mm, -26mm)  {Str "then side"};
  \node[ast off] (s3)   at (74mm, -26mm)  {Str "else side"};
  \node[ast off] (vn)   at (-34mm, -44mm) {Var n};
  \node[ast off] (n2)   at (-2mm, -44mm)  {Num 2};

  \draw[flow] (call.south) -- (s1.north);
  \draw[flow] (call.south) -- (cond.north);
  \foreach \c in {lt, s2, s3} { \draw[flow, color=TcGrayText] (cond.south) -- (\c.north); }
  \draw[flow, color=TcGrayText] (lt.south) -- (vn.north);
  \draw[flow, color=TcGrayText] (lt.south) -- (n2.north);

  % ---- コマ10 の到達範囲 ----
  \begin{scope}[on background layer]
    \node[fill=TcAccentLighter, rounded corners=2mm, fit=(call)(s1), inner sep=3mm] (reach) {};
  \end{scope}
  \node[note, anchor=east] at ($(reach.west)+(-2mm,0)$)
      {コマ10 の {\ttfamily collect\_strings\_expr()} が\\
       枝を持つのは {\ttfamily Str} / {\ttfamily Assign} / {\ttfamily Call} の\\
       3 種類だけ（実線の範囲）};
  \node[note, anchor=west] at ($(cond.east)+(4mm,0)$)
      {{\ttfamily Cond} はコマ10 では {\ttfamily case \_: pass} に落ちる。\\
       ここで走査が止まり、破線の枝は素通りしていた};

  % ---- コマ11 の走査順 ----
  \node[num, anchor=south east] at ($(call.north west)+(-0.5mm,0.5mm)$) {①};
  \node[num, anchor=south east] at ($(s1.north west)+(-0.5mm,0.5mm)$)   {②};
  \node[num, anchor=south east] at ($(cond.north west)+(-0.5mm,0.5mm)$) {③};
  \node[num, anchor=south east] at ($(lt.north west)+(-0.5mm,0.5mm)$)   {④};
  \node[num, anchor=south east] at ($(s2.north west)+(-0.5mm,0.5mm)$)   {⑤};
  \node[num, anchor=south east] at ($(s3.north west)+(-0.5mm,0.5mm)$)   {⑥};
  \node[note, anchor=west] at (18mm, -45mm)
      {①〜⑥ はコマ11 の走査がたどる順番。{\ttfamily n} と {\ttfamily 2} には\\
       文字列が無いので、たどっても何も登録されない};

  % ---- 集めた結果 ----
  \node[lab, anchor=north] at ($(s1.south)+(0,-1.5mm)$) {.LC1};
  \node[lab, anchor=north] at ($(s2.south)+(0,-1.5mm)$) {.LC2};
  \node[lab, anchor=north] at ($(s3.south)+(0,-1.5mm)$) {.LC3};

  \node[region, minimum width=52mm, minimum height=30mm, anchor=north west] (data) at (116mm, 4mm) {};
  \node[lcline, anchor=north west] at ($(data.north west)+(3mm,-3mm)$) {.data};
  \node[lcline] at ($(data.north west)+(3mm,-12mm)$) {.LC1:  "\%s{\textbackslash}n"};
  \node[lcline] at ($(data.north west)+(3mm,-19mm)$) {.LC2:  "then side"};
  \node[lcline] at ($(data.north west)+(3mm,-26mm)$) {.LC3:  "else side"};

  \draw[ptr] (104mm, -20mm) -- (114mm, -20mm);
  \node[gannot, anchor=north, align=center] at (142mm, -28mm)
      {集めた順に {\ttfamily .LC1} から割り当て、\\
       各ラベルの下に {\ttfamily .byte} 列を並べる};

  % ---- 補足 ----
  \node[gannot, anchor=north, align=center] at (40mm, -60mm)
      {走査は「実行される経路」ではなく「AST に書かれているもの全部」を対象にする。\\
       実行時に選ばれるのは片方でも、両方のラベルが {\ttfamily .data} に無いとリンクに失敗する};
\end{tikzpicture}
\end{document}
